\cotefile{28}
\scanpage{2}
\uline{Demazure} 5.1.1970 \quad $\underline{\mathrm{Aut}}^{\circ}$

\uline{Introduction}
\diagram{\(\begin{array}{ll|l}
P_1\times P_1 && \mathrm{PGl}_2\times\mathrm{PGl}_2\\
P_2 && \mathrm{PGl}_{3}\deleted{\unsure{(3)|(G_3)}}\\
F_n & n\geq\deleted{\unsure{0|2}}1\quad(F_0=P_1\times P_1)
  & \bigl(\unsure{\mathrm{Gl}(2)|\mathrm{Gl}_2|\mathrm{GL}(2)}\mathbin{\cdot}\Gamma(\unsure{\mathcal O(n)|O(n)})\bigr)/\mu_n\\
 & (F_1=P_2\ \text{éclaté}) &
\end{array}\)}

\uline{Th.} (Enriques) : tout ss-gpe de $\mathrm{Cr}_2$, algébrique lisse connexe, est $\subset$ un des précédents à conjugaison près\ldots.

\uline{Cor.}
\[
\left\{\begin{aligned}
&\text{\deleted{\unsure{gpe}}\inserted{partie}}\ \text{réductive \uline{de type} }A_n\\
&\text{ils sont \uline{définis sur} }Z\\
&\text{sont des groupes d'autom. d'espaces très simples, \uline{proj. et lisses}}\\
&\text{(opérant avec orbites ouvertes)}\\
&\text{tous de rang }2
\end{aligned}\right.
\]

On va généraliser à la dim $n$, sauf difficulté liée aux tores maximaux ; OK si on prend les ss-groupes de \uline{rang maximum}

\uline{Sorites} \quad (Cas $X$ lisse séparé \inserted{\unsure{t.f.p.|t.f.}} sur $S$)

Pseudo-morphismes, pseudo-isomorphismes (rel/$S$)

Pseudo-automorphismes \(\mathrm{Psaut}_S(X),\quad\underline{\mathrm{Psaut}}_S(X)\)
\[
\begin{aligned}
&\underline{\mathrm{Psaut}}_k(K),\qquad
\underset{\sim}{\mathrm{Cr}}_{n,k}=\underline{\mathrm{Psaut}}_k(k[\unsure{t_1,-,t_n|t_1,\ldots,t_n}])\\
&\underline{\mathrm{Psaut}}_k(K)(k)\simeq\mathrm{Aut}_k(K)^{\circ}\\
&\mathrm{Lie}(\underline{\mathrm{Psaut}}_k(K))\simeq\mathrm{Der}_k(K)^{\circ}
\end{aligned}
\]

\uline{Proposition} Section unité de $\underline{\mathrm{Psaut}}_S(X)$ imm. fermée dans chacun des 2 cas

a) $S$ artinien

b) fibres \inserted{\unsure{géomt|géom.}} irréductibles

% archive sheet number: 1
